\documentclass[11pt]{article}
\usepackage[margin=1in]{geometry}
\usepackage{amsmath,amssymb}
\usepackage{booktabs}
\usepackage{array}

\title{How Much Is Domain Knowledge Worth?\\ Pricing Hypothesis-Space Constraints in Bits\\ and Auditing Them on a Dated Corpus}
\author{Ranadhir Ghosh\\
FIS Global, Jacksonville, Florida, USA\\
\texttt{ranadhir.ghosh@fisglobal.com}\\
ORCID 0009-0001-5439-4258}
\date{}

\begin{document}
\maketitle

\begin{abstract}
Claims that domain knowledge will rescue a discovery pipeline are usually
rhetorical; this paper makes them auditable. Every structural constraint on
a hypothesis space is summarised by two numbers --- its pruning power
$\rho$ (the fraction of wrong hypotheses rejected) and soundness $\sigma$
(of true hypotheses retained) --- priced by an accounting lemma at
$\log_2(\sigma/(1-\rho))$ bits --- the constraint's positive likelihood
ratio as a diagnostic test. We audit a dated corpus of 23{,}710 post-cutoff
ChEMBL hypotheses (17{,}143 genuine, 6{,}567 spurious; measured error rate $p = 0.2770$).
(i) Relational priors split: molecule-level
guilt-by-association is at chance (AUCs $0.4929$--$0.4954$) and its filters
are the study's most destructive objects; a target's pre-cutoff hit-rate is
informative (AUC $0.7266$), its tractability filter earning $+0.857$ bits
--- under half the bar. (ii) All five canonical drug-likeness rules are
harmful: negative bits, edge-level 95\% intervals excluding zero for four
of the five (Veber's indistinguishable from zero). (iii) Physicochemical
structure supplies at most $+0.729$ bits under oracle-optimistic
conditions, and $1.206$ bits in a leakage-free temporal holdout adding
target identity, against $1.85$ bits required under the loop model of
Section 3 ($1.99$ bits, static). Under an explicit threshold model, a sound
constraint provably rescues a collapsed loop (positive bits, not pruning
alone, are required --- correcting an earlier statement); an unsound one
caps reachable truths at a fraction $\sigma$. The contribution is the
audit: ``add domain knowledge'' becomes a threshold-comparable number ---
and on this corpus, the canonical knowledge fails it.
\end{abstract}

\noindent\textbf{Keywords:} autonomous scientific discovery; inductive bias;
hypothesis space; phase transition; false discovery; AI for science.

\section{Introduction}

\textbf{Research question.} Is the hypothesis-error rate of a closed
discovery loop a controllable quantity rather than an exogenous one; can a
structural constraint on the hypothesis space be priced in bits of evidence;
and does any constraint available in a real scientific corpus supply enough
bits to move a collapsed loop back across its throughput threshold?

\textbf{Significance.} The companion throughput theory \cite{ghosh-throughput}
gives operators four levers, all of which cost money: buy verification, buy
capacity, cut reuse, or propose less. This paper prices a fifth that costs
nothing at run time --- knowing something true about the problem --- on the
same scale, so that ``we should add domain knowledge'' becomes a number that
can be audited against a threshold. Because the price is explicit, the audit
also convicts false constraints, and we show it doing so on rules that are in
everyday industrial use.

This paper extends the companion throughput theory
\cite{ghosh-throughput} by supplying its missing upstream lever. It is
written to be refereable on its own: Section 3 states every part of that
theory the argument uses, with its operating points, and Section 3.3 lists
explicitly what is derived here and what is imported.

\textbf{Contributions.}

\begin{enumerate}
\item \textbf{An audit protocol and its accounting lemma (Section 4).} A
constraint is a pair $(\rho, \sigma)$; it shifts error log-odds additively by
$\log_2(\sigma/(1-\rho))$. The identity is Bayes' rule in odds form --- the
multiplier is the positive likelihood ratio of the constraint read as a
diagnostic test (Section 2) --- and we make no deeper claim for it. Its value
is operational: it puts every knowledge claim on one scale, against one
threshold, measurable on a dated corpus.
\item \textbf{Measurement on a real corpus, and an honest negative
(Section 9).} On 23{,}710 dated ChEMBL hypotheses the requirement is
$1.85$--$1.99$ bits; molecule-level relational priors supply $\leq 0$
and target tractability $+0.857$; all five canonical drug-likeness
rules are harmful; structure supplies at most $+0.729$ bits even under
oracle conditions, and $1.206$ bits in a leakage-free temporal
holdout when target identity is added. Section 9.5 reports uncertainty intervals for every headline
quantity.
\item \textbf{Threshold consequences under an explicit model (Sections 3,
5--8).} A corrected knowledge-rescue theorem (the guarantee requires
positive bits; pruning alone can raise the error rate, as the corpus itself
demonstrates); a matched-volume control attributing the rescue to
selectivity rather than volume; a ceiling-and-scaling statement of what
unsoundness costs; and the knowledge requirement of a published replication
pipeline, computed under both the model-conditional and the static
accounting.
\end{enumerate}

\section{Related work}

Classical sequential analysis \cite{wald} and active learning \cite{settles}
quantify the cost of labelling within a fixed hypothesis class. Work on
false-discovery control \cite{benjamini} prescribes corrections over a fixed
candidate set, and the anomaly replication literature \cite{hou} measures how
badly that set is contaminated in practice. The companion closed-loop theory
\cite{ghosh-throughput} adds the dynamical account in which the candidate set
grows under reuse and derives the interior threshold $p_c$; it treats $p$ as
given. The degradation of a reused candidate pool also has a published
cousin in adaptive data analysis \cite{dwork}, where adaptively reusing a
holdout destroys its validity unless the reuse is budgeted; the companion
theory's reuse lever is the same phenomenon inside a discovery loop, and we
flag the connection rather than rediscover it.

On the knowledge side, physics-inspired symbolic regression prunes search
using dimensional analysis and symmetry \cite{udrescu}, raising solve rates
substantially, but offers no cost theory and no threshold; the geometric deep
learning programme \cite{bronstein} derives architectures from symmetry
without pricing them. That assumptions are indispensable for generalisation
is classical \cite{mitchell, wolpert}, and bias itself has been treated as a
learnable, priceable object before: Baxter's model of inductive-bias
learning \cite{baxter} bounds the sample savings a well-chosen hypothesis
space buys across related tasks, and Boopathy et al.\ \cite{boopathy}
measure inductive-bias complexity directly in bits for supervised tasks. Our
accounting lemma sits in an older lineage still: read as a diagnostic test,
a constraint's multiplier $\sigma/(1-\rho)$ is its positive likelihood ratio,
and ``conditioning on a passed test multiplies the odds by $\mathrm{LR}_+$''
is the standard odds form of Bayes' rule in diagnostic testing
\cite{deeks}. We therefore claim no novelty for the identity. Reviews of
AI-driven discovery \cite{grosso} state the three ingredients --- inductive
bias is unavoidable, sample complexity bounds learning, experimental
constraints bind --- without a model linking them, while autonomous
laboratories \cite{king, szymanski} report throughput without a predictive
account of its ceiling. Our delta is narrower than a first: to our knowledge
this is the first study to price a structural constraint \emph{in the same
units as a closed-loop discovery-throughput threshold}, and to run that
audit, with leakage controls and uncertainty intervals, on a real dated
corpus (ChEMBL \cite{mendez}). The rules we audit are those of Lipinski et
al.\ \cite{lipinski}, Veber et al.\ \cite{veber}, and Bickerton et al.\
\cite{bickerton}.

\section{The loop and its threshold: the model, stated in full}

Earlier versions of this work left the threshold machinery in the companion
manuscript. Everything the present argument uses is stated here.

\subsection{The loop and the imported threshold law}

A \emph{loop} runs in rounds. Each round the generator proposes $H$
hypotheses, of which a fraction $p$ are wrong (the \emph{hypothesis-error
rate}); a verification stage validates proposals subject to a capacity $V$
(the maximum number of validations per round); a fraction $R$ of each
round's proposal pool is drawn from previously seen candidates (\emph{reuse}).
The loop's output is its \emph{validated-discovery throughput} $\bar D$,
the mean number of validated true discoveries per round. A pipeline is
\emph{productive} if $\bar D > 0$ in the long run and \emph{collapsed} if
$\bar D = 0$.

From the companion throughput theory \cite{ghosh-throughput} we import
exactly four statements, and nothing else:

\begin{itemize}
\item[(M1)] \textbf{Interior threshold.} For fixed $(H, V, R)$ there is a
critical error rate $p_c(H, V, R) \in (0, 1)$ such that the mean-field
throughput is zero for $p > p_c$ and positive for $p < p_c$. ``Zero'' is the
mean-field statement: finite loops simulated near $p_c$ exceed the
mean-field prediction by a small finite-size effect (Section 11), which is
why one arm of Table 1 shows throughput $0.068$ while sitting above its
threshold.
\item[(M2)] $\partial p_c / \partial H < 0$: proposing less raises the
tolerable error rate.
\item[(M3)] $\partial p_c / \partial V > 0$: buying verification raises it.
\item[(M4)] $\partial p_c / \partial R < 0$: cutting reuse raises it.
\end{itemize}

Statements (M2)--(M4) are the companion theory's four priced levers. We write
$p_c$ for the threshold at the operating point in force and $p_c'$ for the
threshold after a constraint has shrunk the stream to $H'$.

\subsection{Operating points and the threshold schedule}

Two studies in this paper share one operating point, and a third uses
another; we tabulate them here so that every downstream number can be
checked.

\emph{Simulation and ChEMBL studies (Sections 5--7 and 9).} At the
unfiltered volume, $p_c = 0.0882$. The companion model's threshold rises as
the retained fraction $r = H'/H$ falls; over the tabulated range the
schedule is, to the precision printed in Tables 1--2,
\begin{equation}
p_c(r) \;=\; 0.0882 \;+\; 0.038\,(1 - r), \qquad r \in [0.800, 1].
\label{eq:schedule}
\end{equation}
This schedule reproduces every $p_c'$ value in Table 1
($0.0890, 0.0896, 0.0904$ at $r = 0.9785, 0.9641, 0.9425$), the rescue
pruning $\rho^\star = 0.2794$ of Section 5, and the ChEMBL bar of Section 9
($\rho^\star = 0.7234$ at $r = 0.800$, where it gives $p_c' = 0.0958$).
The lower endpoint of the stated domain is that Section 9 rescue point
itself: at $r = 0.800$ the schedule implies $\rho^\star = 0.7234$, agreeing
with the printed $0.7234$ to within the rounding of the printed inputs
(the two domain floors differ by $0.00001$ in $p_c$).
% AUTHOR INPUT NEEDED: eq. (2) is the linear schedule that reproduces all
% tabulated p_c' values of the companion model at this operating point to
% printed precision. The author should confirm it against the companion
% manuscript, or replace it with the companion's closed form.

\emph{Replication-anchor study (Section 8).} The published-factor pipeline
of \cite{hou} runs at $p_c = 0.125$ under strong governance and
$p_c = 0.032$ under weak governance, where \emph{governance} names the
pipeline's bundle of verification-capacity and reuse settings $(V, R)$.
That pipeline's volume schedule is not reproduced here; Section 8 therefore
reports its knowledge requirement under two accountings, one of which (the
static accounting) needs no schedule at all.
% AUTHOR INPUT NEEDED: the (H, V, R) values behind the two governance
% regimes of Table 3, and that pipeline's p_c(r) schedule, live only in the
% companion record. Supply them, or Table 3's model-conditional column
% remains labelled as imported.

\subsection{What is derived here, and what is imported}

Derived in this paper: the accounting lemma (Section 4); the corrected
rescue theorem and the ceiling/scaling statements (Sections 5 and 7), which
are proved from the lemma plus (M1)--(M2); every $p'$ and bits figure in
Tables 1--8, each of which is arithmetic on printed inputs; the
static-threshold requirements (Sections 8 and 9), which use the lemma
alone; and the intervals of Section 9.5. Imported from the companion record
\cite{ghosh-throughput}: the threshold law (M1)--(M4) and its operating
points; the schedule (\ref{eq:schedule}); the model-conditional column of
Table 3; and the throughput columns of Tables 1--2 together with the
simulation figure of Section 7, which are outputs of the companion
simulator (Table 1: mean over 12 seeds) reported as such, not derived here.

\section{The constraint operator and the accounting lemma}

The loop proposes $H$ hypotheses per round, of which $H(1-p)$ are correct.

\medskip
\noindent\textbf{Definition 1 (Constraint).} A constraint $C$ applied to the
proposal stream is the pair $(\rho, \sigma)$ where $\rho$ is the fraction of
wrong hypotheses it rejects (its \emph{pruning power}) and $\sigma$ the
fraction of true hypotheses it retains (its \emph{soundness}). $C$ is
\emph{sound} iff $\sigma = 1$.

After filtering,
\begin{equation}
p' = \frac{(1-\rho)\,p}{(1-\rho)\,p + \sigma(1-p)}, \qquad
H' = H\big[\sigma(1-p) + (1-\rho)p\big].
\end{equation}

\medskip
\noindent\textbf{Lemma 1 (Accounting identity).}
$\log_2 \mathrm{odds}(p') = \log_2 \mathrm{odds}(p) -
\log_2 \frac{\sigma}{1-\rho}$. The constraint therefore supplies
$\log_2 \frac{\sigma}{1-\rho}$ \emph{bits} of evidence about hypothesis
correctness, additively and at zero experimental cost.

\medskip
\noindent\emph{Proof.} The retained stream contains $(1-\rho)p$ wrong and
$\sigma(1-p)$ true hypotheses (as fractions of the original stream), so
$\mathrm{odds}(p') = \frac{(1-\rho)p}{\sigma(1-p)} =
\mathrm{odds}(p)\cdot\frac{1-\rho}{\sigma}$. Taking $\log_2$ gives the
identity. \hfill$\square$

\medskip
\noindent\textbf{Lineage.} Read as a diagnostic test \cite{deeks}, the
constraint's retention rate on true hypotheses is a sensitivity ($\sigma$)
and its retention rate on wrong hypotheses a false-positive rate ($1-\rho$);
the multiplier $\sigma/(1-\rho)$ is the test's positive likelihood ratio
$\mathrm{LR}_+$, and Lemma 1 is Bayes' rule in odds form. This is the Bayes
budget in explicit form: evidence from examples and evidence from
assumptions enter the same ledger. We use the lemma as bookkeeping --- the
instrument that makes Sections 5--9 auditable --- not as a characterisation
result.

Note the sign. A constraint with $\sigma < 1-\rho$ supplies \emph{negative}
bits: it removes more truth than error and actively degrades the generator.
Section 9 exhibits five such rules in industrial use, plus a family of
relational filters with the same property --- which is why ``pruning'' alone
can never be the criterion: the sign of the bits is.

\medskip
\noindent\textbf{Remark (composition).} Constraints applied in sequence
compose additively in bits when each constraint's $(\rho, \sigma)$ is
measured on the stream it actually receives: filtering by $C_1$ then $C_2$
multiplies the odds by $\mathrm{LR}_+(C_1) \cdot \mathrm{LR}_+(C_2 \mid
\text{stream after } C_1)$. Pairs $(\rho, \sigma)$ measured independently on
the same unfiltered corpus do \emph{not} compose this way. ``Independence''
in this paper always means the sequential, on-stream sense.

\medskip
\noindent\textbf{Remark (bits price selectivity, not yield).} As
$\rho \to 1$ at fixed $\sigma > 0$, the bit count grows without bound while
the retained stream shrinks to $H\sigma(1-p)$ and the reachable truths stay
capped at a fraction $\sigma$ (Section 7). Bits must therefore always be
read together with $\sigma$ and with $H'$. A filter can prune 99\% of a
stream while retaining 0.19\% of its truths --- Table 4, row 1 --- and the
bit price alone does not say whether that trade serves the loop.

\medskip
\noindent\textbf{Numerical check.} Over 20,000 random $(p, \rho, \sigma)$
triples the maximum absolute discrepancy in Lemma 1 is $3.65\mathrm{e}{-12}$
bits. This verifies the floating-point evaluation of the algebra --- the
identity against itself --- and nothing about any loop or corpus; we report
it so that no reader mistakes it for evidence.

\section{Knowledge rescue}

\noindent\textbf{Proposition 2 (Double benefit, corrected).} Let
$C = (\rho, \sigma)$ act on a loop at $(p, H)$, and write
$b = \log_2 \frac{\sigma}{1-\rho}$ for its bit price.

\begin{itemize}
\item[(i)] $p' < p$ iff $b > 0$; $p' = p$ iff $b = 0$; $p' > p$ iff $b < 0$.
\item[(ii)] $H' < H$ for every non-identity constraint ($\rho > 0$ or
$\sigma < 1$); hence, by (M2), $p_c' \geq p_c$, strictly under the schedule
(\ref{eq:schedule}) whenever $H' < H$.
\item[(iii)] \emph{Rescue.} If $C$ is sound ($\sigma = 1$) and the loop is
collapsed ($p > p_c$), there exists $\rho^\star < 1$ such that every
$\rho > \rho^\star$ renders the loop productive.
\item[(iv)] For an unsound constraint with $b > 0$, both sides of the
productivity condition move favourably ($p$ falls, $p_c$ rises) but no
$\rho^\star < 1$ is guaranteed in general; and for $b < 0$ the two sides
move in \emph{opposite} senses --- $p$ rises while $p_c$ rises --- so a
pruning constraint can deepen a collapse. Table 4's first row is the corpus
instance: $\rho = 0.9889$ raises $p$ from $0.2770$ to $0.6952$.
\end{itemize}

\noindent\emph{Proof.} (i) By Lemma 1, filtering multiplies
$\mathrm{odds}(p)$ by $(1-\rho)/\sigma = 2^{-b}$, and $p'$ is a strictly
increasing function of its odds. (ii) $H'/H = \sigma(1-p) + (1-\rho)p
= 1 - \big[(1-\sigma)(1-p) + \rho p\big] < 1$ unless $\sigma = 1$ and
$\rho = 0$. (iii) For $\sigma = 1$, $p'(\rho)$ is continuous and strictly
decreasing with $p'(\rho) \to 0$ as $\rho \to 1$, while
$H'(\rho) \to H(1-p) > 0$ and $p_c(H'(\rho))$ stays at or below
$p_c(H(1-p))$ (since $H'(\rho) \geq H(1-p)$ and $p_c$ is decreasing in
volume) yet at or above $p_c(H) > 0$ (since $H'(\rho) \leq H$); hence
$g(\rho) = p_c(H'(\rho)) - p'(\rho)$ is continuous,
negative at $\rho = 0$ (the loop is collapsed), and positive near
$\rho = 1$, so a crossing $\rho^\star$ exists, and $g$ is strictly
increasing, so the crossing is unique. \hfill$\square$

Two qualifications are part of the statement, not footnotes to it. First,
the guarantee in (iii) is proved for sound constraints; the original,
stronger slogan --- any $\rho > 0$ reduces $p$ --- is false, and this
paper's own Table 4 is the counterexample (iv). Second, ``productive''
certifies the \emph{sign} of $p_c' - p'$: positive long-run throughput in
the model of Section 3. It does not by itself lower-bound the discovery
\emph{rate}, which also scales with the retained volume $H'$; a
high-$\rho$ filter can satisfy the condition while discovering very little
per round. The economics of rescue --- that knowledge is ``an independent,
and often cheaper, lever'' --- are therefore a claim about the threshold,
priced in Section 8 in the units an operator actually faces.

Table 1 shows a loop moving from permanent zero throughput to 33.284
validated discoveries per round. No additional experiments were purchased;
the entire change is a restriction of the hypothesis space.

\begin{table}[t]
\centering
\caption{Rescue of a supercritical loop by a sound constraint. Operating
point $p = 0.12$ against $p_c = 0.0882$; $\rho^\star = 0.2794$ under the
schedule (\ref{eq:schedule}), at which $p' = p_c' = 0.0895$. Throughput is
the simulated mean over 12 seeds (companion simulator, Section 3.3). The
$\rho = 0.179$ arm sits above its shifted threshold yet shows throughput
$0.068$, not $0$: the finite-size effect of (M1), cf.\ Section 11.}
\begin{tabular}{lccccc}
\toprule
arm & $\rho$ & $p'$ & $p_c'$ & productive & throughput \\
\midrule
none & 0.000 & 0.1200 & 0.0882 & no & 0.000 \\
sound below threshold & 0.179 & 0.1006 & 0.0890 & no & 0.068 \\
sound at threshold & 0.299 & 0.0872 & 0.0896 & yes & 7.845 \\
sound strong & 0.479 & 0.0663 & 0.0904 & yes & 33.284 \\
\bottomrule
\end{tabular}
\end{table}

\section{The control: selectivity, not volume}

A constraint shrinks the proposal stream, and shrinking $H$ alone raises
$p_c$ (M2). The rescue could therefore be an artefact of reduced volume with
no knowledge in it. We compare against a random filter discarding the same
total fraction of proposals but blind to truth, which by construction leaves
$p$ unchanged.

\begin{table}[t]
\centering
\caption{Matched-volume control. Both arms discard the same fraction of
proposals.}
\begin{tabular}{lcccccc}
\toprule
& & retained & \multicolumn{2}{c}{knowledge filter} &
\multicolumn{2}{c}{random filter} \\
$\rho$ & & & $p'$ & throughput & $p'$ & throughput \\
\midrule
0.299 & & 0.964 & 0.0872 & 7.845 & 0.1200 & 0.000 \\
0.479 & & 0.942 & 0.0663 & 33.284 & 0.1200 & 0.000 \\
0.900 & & 0.892 & 0.0135 & 70.467 & 0.1200 & 0.000 \\
\bottomrule
\end{tabular}
\end{table}

Volume reduction alone rescues nothing at these parameters. The effect is
attributable to selectivity. We note what this control does and does not
show: inside a model where throughput is a function of $(p, H)$, a
truth-blind filter moves only $H$, so the null result attributes the rescue
to selectivity rather than volume \emph{within the model}; it cannot
validate the model against a real loop, and no closed loop is run in this
paper (Section 11).

\section{The cost of unsoundness}

If $\sigma < 1$ the constraint deletes true hypotheses, so the reachable set
is capped. To state the cost precisely we need two quantities the earlier
draft left undefined. Let $T$ be the number of true hypotheses in the
declared universe --- the \emph{attainable truths} --- and let $\bar D_0$ be
the mean validated discoveries per round of the unconstrained pipeline over
the horizon in question.

\medskip
\noindent\textbf{Proposition 3 (Ceilings and overtaking, restated).}
\begin{itemize}
\item[(a)] \emph{Ceiling.} A constrained pipeline can validate at most
$\sigma T$ truths in total; its ceiling equals the unconstrained
pipeline's ($T$) iff $\sigma = 1$.
\item[(b)] \emph{Rescue regime.} If the unconstrained pipeline is collapsed
(throughput $0$) and the constrained one is productive, the constrained
pipeline is never overtaken: the unconstrained pipeline discovers nothing,
at any horizon.
\item[(c)] \emph{Scaling horizon.} If the unconstrained pipeline sustains a
constant mean rate $\bar D_0 > 0$ until it has accumulated $\sigma T$
discoveries --- an assumption, violated by any saturating process near its
own ceiling --- then it draws level with the constrained pipeline's
permanent ceiling after $n \approx \sigma T / \bar D_0$ rounds.
\end{itemize}

\noindent\emph{Proof.} (a) The pipeline can only validate truths the filter
retains; the filter retains a fraction $\sigma$ of true hypotheses, so at
most $\sigma T$ of the universe's truths are ever reachable. (b) Immediate
from the definitions. (c) Arithmetic under the stated constant-rate
assumption. \hfill$\square$

What soundness alone does \emph{not} decide is finite-horizon overtaking in
general: with $\sigma = 1$ the ceilings coincide, but if the unconstrained
loop is productive and faster per round, it reaches any fixed discovery
count first. The earlier headline form of this result --- ``soundness alone
decides whether a constrained pipeline is ever overtaken'' --- is true only
in the rescue regime (b) and as a statement about ceilings (a), and we
restrict it to those. The horizon in (c) is a scaling law, not an absolute
round count (Section 11).

This ceiling law is the same shape as the small-data advantage of
architecturally constrained learners over flexible ones, derived here in an
entirely different formalism. We note one sharp consequence visible in the
companion simulator's computations: an unsound constraint with no pruning
power ($\rho = 0$, $\sigma = 0.5$) supplies exactly $-1$ bit (Lemma 1),
doubling the error odds --- on the simulated loop, which sat near its
threshold, this took an already-productive pipeline from 38.101 to 0.000
discoveries per round. (The direction is forced by the arithmetic; the
figure itself is a reported output of the companion simulator at an
operating point not tabulated in the printed record, cf.\ Section 3.3.) A
false constraint does not merely cost coverage; it can extinguish the
pipeline.

\section{How much knowledge a real pipeline needs}

Applying Lemma 1 to the published-factor replication anchor \cite{hou} gives
the knowledge requirement directly. Table 3 reports it under two
accountings. The \emph{static} columns are derivable in this paper from the
lemma alone: the pruning a sound constraint needs if the threshold does not
move, $(1-\rho^\star) \leq \mathrm{odds}(p_c)/\mathrm{odds}(p)$. The
\emph{model-conditional} columns are the values computed under the companion
model, in which the threshold also rises as the stream shrinks (schedule
imported; Section 3.2). The two accountings bracket the requirement.

\begin{table}[t]
\centering
\footnotesize
\caption{Knowledge required to return the published-factor pipeline to
productivity. Anchor: Hou, Xue \& Zhang (2020), Replicating Anomalies.
Model-conditional $\rho^\star$ and bits are imported from the companion
model at that pipeline's operating points (Section 3.2); static values are
derived here from Lemma 1.}
\begin{tabular}{llccccccc}
\toprule
& & & & \multicolumn{2}{c}{model-conditional} &
\multicolumn{2}{c}{static} \\
governance & anchor & $p$ & $p_c$ & $\rho^\star$ & bits &
$\rho^\star$ & bits \\
\midrule
strong governance & $t \geq 1.96$ & 0.650 & 0.125 & 0.8892 & 3.17 & 0.9231 & 3.70 \\
strong governance & $t \geq 2.78$ & 0.821 & 0.125 & 0.9469 & 4.24 & 0.9689 & 5.00 \\
weak governance & $t \geq 1.96$ & 0.650 & 0.032 & 0.9768 & 5.43 & 0.9822 & 5.81 \\
weak governance & $t \geq 2.78$ & 0.821 & 0.032 & 0.9897 & 6.60 & 0.9928 & 7.12 \\
\bottomrule
\end{tabular}
\end{table}

Three to seven bits of genuine structural knowledge would move that pipeline
back across its own transition, under either accounting --- a modest
requirement by the standards of an exact symmetry, and a concrete, auditable
model-risk control.

\section{Measurement on a real corpus}

We now measure $(\rho, \sigma)$ rather than positing it. The corpus is the
dated ChEMBL \cite{mendez} temporal-holdout landscape: a hypothesis is a
post-cutoff candidate (target, compound) edge, true if the measured potency
clears the threshold and wrong otherwise. The execution is specified
in full in the Reproducibility section: ChEMBL release 36, a
document-year cutoff of 2018, a potency threshold of pChEMBL
$\geq 6.0$, and the eight highest-volume human single-protein targets
under a selection rule fixed before execution; the edge-level dataset
is deposited with the code.
% Design facts (cutoff, threshold, target identities) are stated in
% the Reproducibility section; the study was executed to a design fixed
% before execution (see the deposited pre-registration).

\textbf{Leakage rule.} A constraint may read only information available at
proposal time. Measured potency (the label), post-cutoff study effort, and
the candidate's own first-appearance year are excluded by construction.
Physicochemical descriptors are admissible because they are computed from
the candidate molecule's own structure, which is known before any assay.

\textbf{The bar.} With 23{,}710 hypotheses, 17{,}143 genuine and 6{,}567
spurious, the measured hypothesis-error rate is $p = 0.2770$ against
$p_c = 0.0882$: the corpus, run through the loop at its pre-registered
operating point \cite{ghosh-throughput}, is collapsed. The requirement on a
sound constraint comes in the two accountings of Section 8. Against a static
threshold, Lemma 1 alone requires $(1-\rho^\star) \leq
\mathrm{odds}(p_c)/\mathrm{odds}(p) = 0.2525$, i.e.\ $\rho^\star = 0.7475$,
or $1.986$ bits. Under the loop model the threshold rises as the stream
shrinks: at the rescue point the retained fraction is $r = 0.800$ and the
schedule (\ref{eq:schedule}) gives $p_c' = 0.0958$, relaxing the requirement
to $\rho^\star = 0.7234$, i.e.\ $1.854$ bits. We quote the bar as
$1.85$--$1.99$ bits; Section 9.5 shows the full sensitivity range over
plausible schedules is $1.79$--$1.99$ bits, so no conclusion below turns on
which accounting is used.

\begin{table}[t]
\centering
\caption{Pre-cutoff relational constraints on ChEMBL. Negative bits indicate
a filter that makes the hypothesis-error rate worse. Bits are computed from
unrounded $(\rho, \sigma)$; recomputation from the displayed values can
differ by up to $\approx 0.1$ bits on the extreme rows, where $\sigma$ is
very small (cf.\ Section 9.5); the $\geq$3 row retains no true edge at
all, so its bits are $-\infty$.}
\begin{tabular}{lcccc}
\toprule
filter & $\rho$ & $\sigma$ & $p'$ & bits \\
\midrule
molecule was a potent binder of $\geq$1 target pre-cutoff & 0.9889 & 0.0019 & 0.6952 & -2.574 \\
molecule was a potent binder of $\geq$2 targets pre-cutoff & 0.9988 & 0.0001 & 0.8000 & -3.384 \\
molecule was a potent binder of $\geq$3 targets pre-cutoff & 0.9997 & 0.0000 & 1.0000 & $-\infty$ \\
molecule seen at all pre-cutoff & 0.9837 & 0.0022 & 0.7431 & -2.916 \\
target hit-rate pre-cutoff $\geq$ 0.5 & 0.4713 & 0.9574 & 0.1746 & 0.857 \\
target hit-rate pre-cutoff $\geq$ 0.7 & 0.4992 & 0.9115 & 0.1739 & 0.864 \\
both: potent molecule AND tractable target & 0.9963 & 0.0008 & 0.6316 & -2.162 \\
\bottomrule
\end{tabular}
\end{table}

\subsection{Relational priors: no signal from molecules, some from targets}

Not one molecule-side prior supplies positive evidence; the target side is
different. The rank AUC of the underlying features explains the split: mol
potent pre 0.4954, mol seen pre 0.4929, tgt potent pre
0.6689, tgt hitrate pre 0.7266 (0.5 denotes no signal). The
molecule-level guilt-by-association prior is indistinguishable from
chance. This corroborates, by an independent route, the companion
criterion paper's \cite{ghosh-criterion} finding that the same cheap
prior underperforms random selection on this corpus. A target's
pre-cutoff hit-rate, by contrast, is the strongest single cheap feature
measured, and the tractability rows of Table 4 are its operating points:
selective because exactly one of the eight targets --- the hERG
counterscreen channel, pre-cutoff hit-rate 0.2467 --- falls below 0.5,
and worth $+0.864$ bits at the 0.7 cut. Section 9.5 adds the
intervals: the molecule-level AUCs sit within noise of $0.5$, the
extreme molecule rows rest on very small retained counts, and the
tractability rows' edge-level intervals exclude zero --- though their
by-target bootstrap, over only 8 clusters, does not: the tractability
gain is one target deep.

\subsection{Structure is informative but insufficient}

Every structural descriptor supplies positive bits at its best cut. But
the best attainable value, $+0.729$ bits with hindsight-chosen
thresholds and an in-sample multivariate fit, still falls short of the
model-conditional bar by $\approx 1.12$ bits, and of the static bar by
$1.256$ bits. Section 9.5 shows its interval, $[0.65, 0.80]$, lies entirely
below even the lowest plausible bar: the shortfall is not a sampling artefact.

\begin{table}[t]
\centering
\caption{Physicochemical descriptors, 23{,}675 edges with resolved structure.
Thresholds are chosen with hindsight subject to $\sigma \geq 0.50$, so these
are optimistic upper bounds.}
\begin{tabular}{lccccc}
\toprule
descriptor & AUC & $\rho$ & $\sigma$ & bits \\
\midrule
MolWt & 0.6234 & 0.6679 & 0.5006 & 0.592 \\
MolLogP & 0.5401 & 0.5512 & 0.5188 & 0.209 \\
TPSA & 0.6295 & 0.6906 & 0.5017 & 0.697 \\
NumHDonors & 0.5700 & 0.4921 & 0.5992 & 0.238 \\
NumHAcceptors & 0.6344 & 0.5279 & 0.6856 & 0.538 \\
NumRotatableBonds & 0.5625 & 0.4928 & 0.6167 & 0.282 \\
RingCount & 0.5991 & 0.6055 & 0.5446 & 0.465 \\
NumAromaticRings & 0.5639 & 0.3365 & 0.7719 & 0.218 \\
FractionCSP3 & 0.4998 & 0.3378 & 0.6922 & 0.064 \\
HeavyAtomCount & 0.6262 & 0.6138 & 0.5525 & 0.516 \\
QED & 0.4064 & 0.6302 & 0.5015 & 0.440 \\
LipinskiViolations & 0.5687 & 0.0012 & 0.9995 & 0.001 \\
all descriptors (oracle, in-sample) & 0.6515 & 0.6982 & 0.5003 & 0.729 \\
\bottomrule
\end{tabular}
\end{table}

One row of Table 5 invites misreading and we flag it: the
LipinskiViolations \emph{descriptor}, swept as a score under the constraint
$\sigma \geq 0.50$, has its hindsight-optimal cut at the near-trivial one ---
retain every edge with at most three violations, $+0.001$ bits. Table 6 instead applies the canonical
\emph{rule} itself as the filter, which forces the cut and loses $0.316$
bits. The two rows measure different objects: what the violation count is
worth as a score (a thousandth of a bit, at its best cut) versus what the rule costs as a
mandate (negative).

\subsection{The canonical drug-likeness rules are false constraints}

All five standard filters supply negative bits on this task. The mechanism
is visible in the descriptor table: QED has AUC 0.4064, i.e.\ inverted ---
potency in this corpus rises with molecular size and complexity, so
drug-likeness rules preferentially retain the less potent candidates.
Lipinski's Rule of Five \cite{lipinski} is a true constraint for oral
bioavailability and a false one for binding potency. The same restriction
can be sound for one question and harmful for another, and the bit
accounting of Lemma 1 is what makes the difference measurable rather than a
matter of opinion. Section 9.5 shows four of the five verdicts are statistically
decisive at edge level; Veber's interval includes zero, so that rule
is harmful in point estimate only.

\begin{table}[t]
\centering
\caption{Named, pre-registered medicinal-chemistry rules, applied without
tuning.}
\begin{tabular}{lcccc}
\toprule
rule & $\rho$ & $\sigma$ & bits & verdict \\
\midrule
Lipinski Ro5 pass (0 violations) & 0.3140 & 0.5510 & -0.316 & harmful \\
MW $\leq$ 500 & 0.2220 & 0.6214 & -0.324 & harmful \\
QED $\geq$ 0.5 (drug-like) & 0.4265 & 0.4393 & -0.385 & harmful \\
Veber (RotB $\leq$ 10 and TPSA $\leq$ 140) & 0.1119 & 0.8870 & -0.002 & harmful \\
lead-like (MW $\leq$ 350 and LogP $\leq$ 3) & 0.9239 & 0.0275 & -1.467 & harmful \\
\bottomrule
\end{tabular}
\end{table}

\subsection{A leakage-free temporal holdout}

Finally we remove both forms of optimism: the structure--potency model is
fitted on pre-cutoff edges only, and its decision threshold is chosen on
pre-cutoff data only.

\begin{table}[t]
\centering
\caption{Temporal holdout: trained on 55{,}525 pre-cutoff edges with
resolved structure, applied to 23{,}675 post-cutoff candidates with
resolved structure.}
\begin{tabular}{lccccccc}
\toprule
arm & AUC train & AUC test & $\rho$ & $\sigma$ & bits (honest) & bits (oracle) \\
\midrule
structure only & 0.5967 & 0.6041 & 0.5467 & 0.5895 & 0.379 & 0.418 \\
structure + target & 0.7415 & 0.7681 & 0.7113 & 0.6659 & 1.206 & 1.318 \\
\bottomrule
\end{tabular}
\end{table}

The best honest value is 1.206 bits, against the oracle bound of 1.318 and a
requirement of 1.85--1.99 bits. 0 of 2 honest arms clear the bar.

\textbf{Where the optimism was.} The three measurements decompose the gap,
and the two arms decompose it differently. For structure alone: fitting
in-sample gave $+0.729$ bits; refitting on the pre-cutoff era and
keeping a hindsight threshold gives 0.418; choosing the threshold
honestly as well gives 0.379. The loss is again mostly temporal
transfer, 0.311 bits, against 0.039 for threshold selection; the
structure-only model's discrimination is weak in both eras and roughly
stable across the cutoff (AUC 0.5967 train, 0.6041 test). For
structure + target the decomposition inverts: the holdout oracle,
1.318 bits, \emph{exceeds} the in-sample oracle, because target
identity is more informative in the test era than in the training era
(AUC 0.7415 train, 0.7681 test) --- plausibly because the panel's
post-cutoff volume concentrates on its most tractable target (JAK2:
6{,}549 of 23{,}675 test edges, pre-cutoff hit-rate
0.8917). Honest threshold choice costs that arm 0.112 bits. A
constraint is only worth the bits it carries forward: the best honest
constraint carries 65\% of what the loop requires (against the
model-conditional bar; 61\% against the static bar), and the
structure-only constraint 20\% (19\% against the static bar).

\subsection{Uncertainty}

Every interval below is computed from this study's own $2 \times 2$ counts
on the stated basis: 23{,}710 post-cutoff edges for the relational rows
(17{,}143 true, 6{,}567 wrong) and 23{,}675 structure-resolved edges
for the structure tables (17{,}117 true, 6{,}558 wrong). \emph{Method.} We
compute Wilson score intervals for $\rho$ and $\sigma$ separately and
propagate them to bits by the endpoint transformation
$[\log_2(\sigma_{lo}/(1-\rho_{lo})), \log_2(\sigma_{hi}/(1-\rho_{hi}))]$ --- a
conservative rectangle. AUC intervals use the Hanley--McNeil standard
error. All edge-level intervals treat edges as independent; because the
edge-level dataset is deposited with the code, the clustering correction
--- a cluster bootstrap by molecule and by target --- is computed and
reported below rather than deferred.

\begin{table}[t]
\centering
\caption{Edge-level 95\% intervals for the headline quantities, computed
from the study's deposited counts. All treat edges as independent and
are lower bounds on the true uncertainty: edges cluster on 8 targets and
on repeated molecules; the cluster bootstrap is reported in the text.}
\begin{tabular}{lccc}
\toprule
quantity & estimate & 95\% interval & basis \\
\midrule
error rate $p$ & 0.2770 & $[0.271, 0.283]$ & counts \\
Ro5 & -0.316 & $[-0.359, -0.273]$ & counts \\
MW $\leq$ 500 & -0.324 & $[-0.359, -0.288]$ & counts \\
QED $\geq$ 0.5 & -0.385 & $[-0.439, -0.330]$ & counts \\
Veber & -0.002 & $[-0.022, +0.019]$ & counts \\
lead-like & -1.467 & $[-1.717, -1.217]$ & counts \\
potent binder $\geq$1 target (Table 4) & -2.574 & $[-3.399, -1.749]$ & counts \\
potent binder $\geq$3 targets (Table 4) & $-\infty$ & $[-\infty, +1.424]$ & counts \\
target hit-rate $\geq$ 0.7 (Table 4) & +0.864 & $[+0.823, +0.906]$ & counts \\
all descriptors, oracle (Table 5) & +0.729 & $[+0.655, +0.804]$ & counts \\
holdout structure + target, honest & +1.206 & $[+1.136, +1.276]$ & counts \\
holdout structure only, honest & +0.379 & $[+0.323, +0.436]$ & counts \\
AUC, mol seen pre (Section 9.1) & 0.4929 & $[0.485, 0.501]$ & Hanley--McNeil \\
\bottomrule
\end{tabular}
\end{table}

Three conclusions survive the intervals. First, four of the five
canonical rules are decisively harmful at edge level: their intervals
exclude zero, by margins of $0.27$--$1.22$ bits. The fifth, Veber's
rule, is indistinguishable from zero ($-0.002$; interval
$[-0.022, +0.019]$): harmful in point estimate, unproven in either
direction. Second, the structure shortfall is not a sampling artefact:
the oracle interval $[0.65, 0.80]$ lies entirely below
the lowest plausible bar. Third, the relational verdict splits by side.
The molecule rows are decisive at edge level --- the $\geq$1-target
filter's interval $[-3.40, -1.75]$ excludes zero by a wide
margin, on retained counts of 32 true and 73 wrong edges --- and
the molecule-level AUC intervals bracket chance. The tractability
row's edge-level interval, $[+0.82, +0.91]$, excludes zero on
the positive side; the bootstrap paragraph below qualifies it.

\textbf{Bar sensitivity.} The model-conditional bar depends on the schedule
slope $\kappa$ of (\ref{eq:schedule}). Solving the rescue condition across
$\kappa$: $\kappa = 0$ (static threshold) gives 1.986 bits;
$\kappa = 0.019$ gives 1.917; $\kappa = 0.038$ (the operating point)
gives 1.854; $\kappa = 0.057$ gives 1.795. The requirement is therefore
$1.79$--$1.99$ bits under any schedule in a $\pm 50\%$ band around the
operating point --- comfortably above the oracle supply under every entry.

\textbf{Clustering.} The 23{,}710 edges cluster on only 8 targets,
and molecules recur across edges; target-level features take at most 8
distinct values. The intervals above are therefore anti-conservative
for exactly the quantities a referee will ask about. With the
edge-level dataset deposited, the cluster bootstrap is computed here
rather than deferred (2,000 resamples, percentile intervals). By
molecule, the intervals track the edge-level ones closely (oracle
$[+0.67, +0.79]$; holdout structure + target
$[+1.15, +1.27]$; the potent-binder $\geq$1 row
$[-3.27, -1.99]$), because few molecules recur across many edges.
By target --- only 8 clusters --- the intervals widen sharply: among
the rules, Ro5, MW $\leq$ 500 and lead-like still exclude zero, but
QED's ($[-0.89, +0.08]$) and Veber's do not; the oracle's is
$[+0.39, +0.91]$; and the tractability row's is
$[-0.02, +2.21]$, which includes zero: with eight targets, whether the
tractability filter helps depends on which targets the panel happens
to contain --- here, on a single counterscreen channel (hERG) sitting
below the tractability line. The $\geq$3-target row retains no true
edge, so its bits are $-\infty$ and no bootstrap interval exists:
every resample is degenerate in the same way.

\section{Discussion}

The negative result is the useful one. The accounting frames a sharp
expectation --- the bits must come from genuine structural knowledge rather
than from the knowledge graph's topology --- and the measurement sharpens
it: molecular topology supplies nothing, target-level tractability
supplies a little, the standard rules supply less than nothing, and
two-dimensional structure supplies roughly two fifths of what is
required even under oracle conditions; the best honest learned
constraint, which adds target identity to structure, supplies roughly
two thirds. The residual is a concrete target
--- binding-site geometry, docking, or a learned structure--activity model
--- rather than a vague appeal to domain expertise.

More broadly, the audit turns a rhetorical move into a procedure. Any
proposed restriction of a hypothesis space can be assigned a $(\rho,
\sigma)$ on a dated corpus and compared against the threshold its pipeline
must clear. Restrictions that fail the test are not merely unhelpful;
Proposition 3 and Lemma 1 together show the two ways they can be ruinous
--- capping the reachable truths at a fraction $\sigma$, and, when their
bits are negative, raising the error rate they were meant to lower.

\section{Caveats}

\begin{itemize}
\item The corpus covers 8 high-volume human single-protein targets ---
three kinases (VEGFR2, EGFR, JAK2), two ion channels (hERG, Nav1.7),
two GPCRs (cannabinoid receptor 2, dopamine D2) and an aspartyl
protease (BACE1) --- selected by the pre-registered volume rule of the
Reproducibility section; target-level features therefore take few
distinct values, and target-level conclusions are specific to this
panel.
\item $(\rho, \sigma)$ are measured on a realised landscape standing in for
a generator's proposal distribution.
\item The companion corpus built from Open Targets is degenerate for this
question: its post-cutoff pool contains 0 spurious edges, so $p = 0$ and
$\rho$ is undefined. It is reported as a limitation, not as a result.
\item Near $p_c$ the simulation exceeds the mean-field prediction (the
finite-size effect documented in the companion throughput paper
\cite{ghosh-throughput}); the two are reported side by side rather than
reconciled.
\item Overtaking horizons scale with a declared universe size and are
therefore a scaling law, not absolute round counts.
\item The edge-level intervals of Section 9.5 treat edges as
independent; the cluster bootstrap reported there is the corrective,
and its by-target intervals, over only 8 clusters, remain wide.
\item The threshold schedule (\ref{eq:schedule}) is the linear form that
reproduces the companion model's tabulated values at the shared operating
point to printed precision; when the companion theory is available in
citable form, its closed form should replace it.
\item Table 3's model-conditional requirements inherit the replication
pipeline's operating points, which are not derivable from this manuscript;
the static columns are the in-paper check on them.
\end{itemize}

\section*{Reproducibility}

Every number in Section 9 regenerates from the rebuild package
deposited with this manuscript: \texttt{build\_corpus.py} (the corpus,
from the pinned ChEMBL 36 extract deposited alongside it),
\texttt{descriptors.py} (RDKit descriptors), \texttt{measure\_relational.py}
(Table 4 and the feature AUCs), \texttt{measure\_structure.py} (Tables 5--6),
\texttt{holdout.py} (Table 7) and \texttt{intervals.py} (the bar, Table 8 and
the cluster bootstrap), sharing \texttt{metrics.py} and \texttt{bar.py};
\texttt{validate.py} recomputes every printed Section 9 number from the
deposited artifacts (edge lists, descriptors, model scores) and diffs
them against this manuscript. The design was fixed before execution in
a pre-registration deposited with the package: ChEMBL release 36;
document-year cutoff 2018, an edge's era being its (target, compound)
pair's first-appearance year and pre-cutoff potency the maximum over
pre-cutoff activities only; potency threshold pChEMBL $\geq 6.0$; and, as
targets, the eight human single-protein targets with the most
pre-cutoff edges (ties broken by post-cutoff edge count, then ChEMBL
ID): CHEMBL4822 (Beta-secretase 1), CHEMBL279 (Vascular endothelial growth factor receptor 2), CHEMBL240 (Voltage-gated inwardly rectifying potassium channel KCNH2), CHEMBL4296 (Sodium channel protein type 9 subunit alpha), CHEMBL253 (Cannabinoid receptor 2), CHEMBL217 (D(2) dopamine receptor), CHEMBL203 (Epidermal growth factor receptor), CHEMBL2971 (Tyrosine-protein kinase JAK2). ``Tractable target'' in Table 4 means pre-cutoff
hit-rate $\geq 0.5$. The numbers in Sections 3--8 are arithmetic on the
inputs stated there or imported from the companion manuscripts, as
Section 3.2 details; the simulation reuses the companion paper's
simulator rather than reimplementing it.

\textbf{Reproducibility gaps remaining.} The following items concern
the companion manuscripts and remain open; this manuscript does not
guess them: (1) citable forms (titles, preprint identifiers) for the
two companion manuscripts \cite{ghosh-throughput, ghosh-criterion};
(2) the $(H, V, R)$ values behind Table 3's two governance regimes and
that pipeline's threshold schedule; (3) the
simulator operating point behind the 38.101 figure of Section 7.

\section*{Statements and Declarations}

\subsection*{Funding}

No external funding.

\subsection*{Competing interests}

The author is employed by FIS Global. The work is methodological and does
not relate to any FIS product.

\subsection*{Data availability}

The study's data derive from ChEMBL release 36, a public resource
(EBI), pinned by the release checksums recorded in the deposited
package. The package --- extraction and measurement code, the
extracted activity and target tables, edge lists, descriptors and
model scores --- accompanies this manuscript and regenerates every
Section 9 number offline.

\subsection*{Author contributions}

R.G.: conception, formal analysis, software, experiments, writing.

\begin{thebibliography}{99}

\bibitem{benjamini} Benjamini, Y. and Hochberg, Y. (1995). Controlling the
false discovery rate. \emph{J. R. Stat. Soc. B} 57(1), 289--300.

\bibitem{baxter} Baxter, J. (2000). A model of inductive bias learning.
\emph{Journal of Artificial Intelligence Research} 12, 149--198.

\bibitem{bickerton} Bickerton, G. R., Paolini, G. V., Besnard, J., Muresan,
S. and Hopkins, A. L. (2012). Quantifying the chemical beauty of drugs.
\emph{Nature Chemistry} 4, 90--98.

\bibitem{boopathy} Boopathy, A., Liu, K., Hwang, J., Ge, S.,
Mohammedsaleh, A. and Fiete, I. R. (2023). Model-agnostic measure of
generalization difficulty. \emph{Proceedings of the 40th International
Conference on Machine Learning}, PMLR 202.

\bibitem{bronstein} Bronstein, M., Bruna, J., Cohen, T. and
Veli\v{c}kovi\'{c}, P. (2021). Geometric deep learning: grids, groups,
graphs, geodesics, and gauges. arXiv:2104.13478.

\bibitem{deeks} Deeks, J. J. and Altman, D. G. (2004). Diagnostic tests 4:
likelihood ratios. \emph{BMJ} 329, 168--169.

\bibitem{dwork} Dwork, C., Feldman, V., Hardt, M., Pitassi, T., Reingold, O.
and Roth, A. (2015). The reusable holdout: preserving validity in adaptive
data analysis. \emph{Science} 349(6248), 636--638.

\bibitem{grosso} Grosso, G., Mikuni, V. and Heinrich, L. (2026). Are we
ready for AI-driven discovery? AI verification before the next fundamental
physics breakthrough. arXiv:2607.10039.

\bibitem{hou} Hou, K., Xue, C. and Zhang, L. (2020). Replicating anomalies.
\emph{Review of Financial Studies} 33(5), 2019--2133.

\bibitem{king} King, R. D. et al.\ (2009). The automation of science.
\emph{Science} 324, 85--89.

\bibitem{lipinski} Lipinski, C. A., Lombardo, F., Dominy, B. W. and Feeney,
P. J. (2001). Experimental and computational approaches to estimate
solubility and permeability in drug discovery and development settings.
\emph{Advanced Drug Delivery Reviews} 46(1--3), 3--26.

\bibitem{mendez} Mendez, D. et al.\ (2019). ChEMBL: towards direct
deposition of bioassay data. \emph{Nucleic Acids Research} 47(D1),
D930--D940.

\bibitem{mitchell} Mitchell, T. (1980). The need for biases in learning
generalizations. Rutgers CBM-TR-117.

\bibitem{settles} Settles, B. (2009). Active learning literature survey.
Univ.\ of Wisconsin--Madison CS-TR-1648.

\bibitem{szymanski} Szymanski, N. J. et al.\ (2023). An autonomous
laboratory for the accelerated synthesis of novel materials. \emph{Nature}
624, 86--91.

\bibitem{udrescu} Udrescu, S.-M. and Tegmark, M. (2020). AI Feynman: a
physics-inspired method for symbolic regression. \emph{Science Advances} 6,
eaay2631.

\bibitem{veber} Veber, D. F., Johnson, S. R., Cheng, H.-Y., Smith, B. R.,
Ward, K. W. and Kopple, K. D. (2002). Molecular properties that influence
the oral bioavailability of drug candidates. \emph{Journal of Medicinal
Chemistry} 45, 2615--2623.

\bibitem{wald} Wald, A. (1945). Sequential tests of statistical hypotheses.
\emph{Ann. Math. Statist.} 16(2), 117--186.

\bibitem{wolpert} Wolpert, D. (1996). The lack of a priori distinctions
between learning algorithms. \emph{Neural Computation} 8(7), 1341--1390.

\bibitem{ghosh-throughput} Ghosh, R. (2026). [Closed-loop discovery
throughput: a threshold theory.] Companion manuscript, unpublished.

\bibitem{ghosh-criterion} Ghosh, R. (2026). [A criterion study of relational
priors on the ChEMBL landscape.] Companion manuscript, unpublished.

\end{thebibliography}

\end{document}
